%=======================================================================
%  CHAPTER 5 — STRUCTURED KNOWLEDGE REPRESENTATION  (4 hrs)
%=======================================================================
\section{Structured Knowledge Representation}

{\color{sub}\footnotesize 4 hrs \gap$\cdot$\gap asked in \mk{38 of 41 papers}
\gap$\cdot$\gap 5 syllabus sub-topics \gap$\cdot$\gap 319 marks, 9.4\,\% of every mark
ever set \gap$\cdot$\gap the 11 sentence sets and the net-to-frame conversion are worked
in \emph{Ch5 -- Structured Knowledge Representation -- Solved Problems}.}

\lead{Read this first:}
\begin{itemize}
  \item \mk{The chapter is two objects and one comparison.} A
        \textbf{semantic net}, a \textbf{frame}, and \textbf{net Vs frame}. Learn those
        three and you have answered nine tenths of what this chapter has ever set.
  \item \mk{The most-asked question is a drawing}:
        \emph{``convert the given sentences into a semantic network''}, in
        \textbf{12 papers} across \textbf{11 sentence sets}, for 5 to 7 marks each, and
        \mk{six of the eleven are the same three figures with the names
        changed}. Answer with a drawing, never with prose. All worked in the companion.
  \item \mk{Second-most-asked, and the safest marks here}:
        \emph{``explain semantic net and frame with examples, list their advantages and
        limitations''} \tS{13}. One comparison table plus two small drawings.
  \item Three reliable short answers, each a fixed list: \textbf{issues in KR} \tF{7},
        \textbf{approaches to KR} \tF{6}, \textbf{how a KR is evaluated} (four adequacies).
  \item \mk{CD and scripts are only 2 papers, but both times they carry
        the whole 8 marks} (\yr{\textbf{71 Bh}} \m{2+5}, \yr{\textbf{77 Ch}} \m{3+5}).
  \item Pairs the examiner reliably asks together: \textbf{issues in KR $+$ draw the
        semantic network}, \textbf{approaches $+$ semantic nets with an example}, and
        \textbf{compare net and frame $+$ convert one into the other}.
\end{itemize}

\hr

%-----------------------------------------------------------------------
\T{5.1 Knowledge, representation, and what makes a representation ``structured''}

\Q{What is knowledge representation? What do you understand by structured knowledge representation?}{\m{2}\m{2+5}}
{\color{sub}\footnotesize \yr{\textbf{75 Bh} \m{2}, \textbf{\textit{76 Bh}} \m{2+5}}
\gap$\cdot$\gap 2 papers at 2 marks each, so no tier chip \gap$\cdot$\gap
\mk{always the opener} to a semantic-net question, so answer it in
three lines and move on}

\sbs{%
\lead{Knowledge representation}
\begin{itemize}
  \item \textbf{Encoding what an agent knows about the world in a form a computer can
        manipulate}, chosen so that \mk{inference is easy}.
  \item Two things exist and must not be confused:
  \begin{itemize}
    \item \textbf{Facts} $-$ truths in the world. The \textbf{knowledge level}.
    \item \textbf{Representations} of those facts in some formalism. The
          \textbf{symbol level}.
  \end{itemize}
  \item KR is the \textbf{mapping} between the two, and it must run \emph{both ways}:
        \emph{Tommy is a dog} $\rightarrow$ \code{dog(Tommy)}; add
        $\forall x: dog(x) \rightarrow hastail(x)$, deduce \code{hastail(Tommy)}; map back
        to \emph{Tommy has a tail}.
  \item \mk{The backward map is where representations fail.} A
        formalism that cannot say its answer back in the user's terms is useless in an
        expert system.
  \item Three \textbf{levels} a fact sits at: \textbf{knowledge} (what is true),
        \textbf{logical} (the sentence encoding it), \textbf{implementation} (the data
        structure holding the sentence).
  \item The \textbf{kinds of knowledge} to name if asked: objects, events, performance
        (how to do a task), meta-knowledge, facts.
\end{itemize}}{%
\lead{Structured knowledge representation}
\begin{itemize}
  \item In logic, everything known about one object is \mk{scattered across
        many separate clauses}. Nothing in the formalism ties them together.
  \item \textbf{Structured KR concentrates all of an object's knowledge in one place}: one
        node of a graph, or one record of slots. The structure itself carries meaning.
  \item So the definition to write: \mk{a representation in which knowledge
        is stored in named, linked units rather than in flat sentences}, each unit holding
        an object and everything known about it.
  \item The four structured schemes on this syllabus: \textbf{semantic nets},
        \textbf{frames}, \textbf{conceptual dependency}, \textbf{scripts}.
  \item Rich and Knight call them \textbf{slot-and-filler} structures, \emph{weak} where
        the slots carry no fixed meaning (nets, frames) and \emph{strong} where they do
        (CD, scripts).
  \item \textbf{What it buys you}: inheritance, default values, and the ability to answer
        ``what do you know about $X$'' by looking in \emph{one} place.
\end{itemize}}

%-----------------------------------------------------------------------
\T{5.2 Approaches, models and schemes of knowledge representation}

\Q{What are the different approaches / models / schemes of knowledge representation? How do you choose among them?}{\m{2}\m{4}\m{7}\m{8}}
{\color{sub}\footnotesize \tF{6} \yr{\textbf{69 Bh} \m{7}, \texttt{80 Ba} \m{2+5},
\textbf{\textit{73 Bh}} \m{8}, 79 Ash \m{4}, \textbf{\texttt{81 Bh}} \m{2},
\textbf{80 Ch} \m{2+2}} \gap$\cdot$\gap
\yr{\textbf{\textit{73 Bh}}} says \emph{``four types of KR scheme''}, which is the
\mk{four-scheme list} below, not the five-method one \gap$\cdot$\gap
\yr{\textbf{80 Ch}} additionally asks \emph{how you make a choice}}

\sbs{%
\lead{The four schemes (answer \yr{\textbf{\textit{73 Bh}}} with these)}
\begin{itemize}
  \item \textbf{1. Logical} $-$ knowledge as formal sentences, manipulated by rules of
        inference.
  \begin{itemize}
    \item Eg propositional logic, first-order predicate logic (Ch\,4).
    \item \checkmark Precise, sound, complete. \xmark Rigid, and a
          single object's facts end up scattered.
  \end{itemize}
  \item \textbf{2. Procedural} $-$ knowledge as procedures that \emph{do} something.
  \begin{itemize}
    \item Eg IF--THEN production rules, an expert system's rule base.
    \item \checkmark Directly executable, control built in.
          \xmark Knowledge and the code using it are inseparable, so
          nothing can be modified independently.
  \end{itemize}
  \item \textbf{3. Network} $-$ knowledge as a labelled graph.
  \begin{itemize}
    \item Eg \textbf{semantic nets}, conceptual dependency, conceptual graphs.
    \item \checkmark Visual, inheritance is free.
          \xmark No quantifiers, no negation.
  \end{itemize}
  \item \textbf{4. Structured} $-$ a network whose nodes are themselves \emph{complex data
        structures}.
  \begin{itemize}
    \item Eg \textbf{frames}, \textbf{scripts}.
    \item \checkmark Defaults, procedures and constraints per slot.
          \xmark No inference engine of its own.
  \end{itemize}
  \item \mk{One sentence that ties them together}: 3 and 4 are the
        \emph{structured} half of this chapter; 1 and 2 are Ch\,4 and the expert system of
        Ch\,7.
\end{itemize}}{%
\lead{The five methods (the wording BJ Sir and PS Sir use)}
\begin{itemize}
  \item \textbf{Procedural} $-$ facts plus the step-by-step operations on them,
        hard-coded.
  \item \textbf{Declarative} $-$ states, facts, rules and relations declared
        \emph{separately} from the algorithm that reasons over them.
  \begin{itemize}
    \item This separation is the whole advantage: the knowledge can be changed without
          rewriting the reasoner, and one optimised reasoner serves many knowledge bases.
  \end{itemize}
  \item \textbf{Relational} $-$ facts as rows of a table, exactly as a database holds
        them. Simple, but carries no inference at all.
  \item \textbf{Hierarchical} $-$ inheritable knowledge as a tree of classes,
        each level inheriting the level above.
  \item \textbf{Complex network graph} $-$ the general labelled graph: semantic nets.
\end{itemize}
\lead{\mk{How to choose} \yr{\textbf{80 Ch}} \m{2}}
\begin{itemize}
  \item Score the candidates on the \textbf{four adequacies} of \S5.3, against the domain
        you actually have.
  \item Then three practical questions:
  \begin{itemize}
    \item Is the knowledge mostly \textbf{taxonomic}? Use a net or frames, and get
          inheritance for nothing.
    \item Does it need \textbf{proof}, quantifiers or negation? Use logic $-$ a net
          cannot express \emph{``no bird is a mammal''}.
    \item Does it need \textbf{defaults and exceptions}? Use frames; logic has no clean
          way to say \emph{``birds fly, except penguins''}.
  \end{itemize}
  \item \mk{The examiner's expected closing line}: no single scheme wins.
        Real systems are \textbf{hybrid} $-$ frames for the objects, rules for the
        reasoning.
\end{itemize}}

%-----------------------------------------------------------------------
\T{5.3 Evaluating a knowledge representation, and the issues in KR}

\Q{How is a knowledge representation evaluated? What are the issues in knowledge representation?}{\m{2}\m{3}\m{4}}
{\color{sub}\footnotesize \tF{7} \yr{\textbf{\texttt{80 Bh}} \m{3}, \textbf{80 Ch} \m{2+2},
79 Ash \m{2+2}, 80 Ash \m{3}, \textbf{\texttt{82 Bh}} \m{2},
\textbf{\textit{76 Bh}} \m{2}, \textbf{\texttt{79 Bh}} \m{2}}
\gap$\cdot$\gap two different lists, and \mk{papers mix them up}:
\emph{properties} is the four adequacies, \emph{issues} is the five design questions
\gap$\cdot$\gap \yr{\textbf{\textit{76 Bh}}} asks the issues inside an expert-system
question}

\sbs{%
\lead{Evaluating a KR: the four adequacies}
\begin{itemize}
  \item \textbf{1. Representational adequacy} $-$ can it represent \emph{all} the kinds of
        knowledge the domain needs?
  \item \textbf{2. Inferential adequacy} $-$ can its structures be manipulated to
        \emph{derive new} structures, ie new knowledge from old?
  \item \textbf{3. Inferential efficiency} $-$ can it carry extra guidance that
        \textbf{focuses} the inference engine in productive directions, rather than
        letting it search blindly?
  \item \textbf{4. Acquisitional efficiency} $-$ can new knowledge be added
        \textbf{easily}, ideally automatically rather than by a human?
  \item Two practical criteria worth adding for the last mark: it must be
        \textbf{easy to use}, and \textbf{easy to modify and extend}.
  \item \mk{The honest closing line}: no known representation scores well
        on all four at once. That is exactly why hybrid systems exist, and it is what
        \yr{\textbf{80 Ch}}'s \emph{``is it possible to evaluate a KR''} is fishing for
        $-$ yes, on these four axes, but only relative to a domain.
\end{itemize}}{%
\lead{The five issues in knowledge representation}
\begin{itemize}
  \item \textbf{1. Important attributes} $-$ are there attributes so basic that they occur
        in almost every domain, and so deserve special handling?
  \begin{itemize}
    \item \mk{Yes: \code{isa} and \code{instance}.} Both carry
          \textbf{property inheritance}, and every scheme in this chapter is built on
          them. This is the expected answer, not a shrug.
  \end{itemize}
  \item \textbf{2. Relationships among attributes} $-$ four to name:
  \begin{itemize}
    \item \textbf{inverses} (\emph{owns} / \emph{owned-by}),
          \textbf{existence in an isa hierarchy} (\emph{height} generalises
          \emph{batting-height}), \textbf{techniques for reasoning about values}
          (types, ranges, rules), \textbf{single-valued attributes} (a person has one
          date of birth).
  \end{itemize}
  \item \textbf{3. Level of representation, ie granularity} $-$ how fine should the
        primitives be?
  \begin{itemize}
    \item Coarse: few facts, but similar sentences look different.
          Fine: identical meanings collapse together, but the store explodes.
          CD (\S5.7) is the extreme fine-grained answer.
  \end{itemize}
  \item \textbf{4. Representing sets of objects} $-$ by \textbf{listing members}, or by
        \textbf{stating the defining property}? Only the second can say
        \emph{``there are lots of sheep''}, which is a property of the set and of no
        member.
  \item \textbf{5. Finding the right structure} $-$ given a huge store, how do you
        \textbf{select} the relevant structure and \textbf{revise} it when it turns out
        wrong? Indexing, matching and frame selection all live here.
  \item Two more the decks add: \textbf{interpretation is informal} in a net, and
        \textbf{complex relationships are hard} to draw with binary links.
\end{itemize}}

%-----------------------------------------------------------------------
\T{5.4 Semantic nets}

\Q{What is a semantic net? Explain with an example. When do we use one?}{\m{2}\m{5}\m{6}\m{7}\m{8}}
{\color{sub}\footnotesize \tS{9} \yr{\textbf{70 Bh} \m{8}, \textit{72 Ma} \m{7},
\texttt{82 Ba} \m{2}, \textbf{75 Bh} \m{6}, \textbf{71 Bh} \m{7}, \textbf{69 Bh} \m{7},
\texttt{80 Ba} \m{5}, \textbf{\textit{76 Bh}} \m{5}, \textbf{\texttt{80 Bh}} \m{5}}
\gap$\cdot$\gap a further \textbf{12 papers} ask you to \emph{draw} one, all worked in
the companion \gap$\cdot$\gap \yr{\textit{72 Ma}} asks for a net
\mk{about the AI course itself}}

\sbsr{0.52}{%
\lead{What it is}
\begin{itemize}
  \item A \textbf{directed labelled graph} that stores knowledge as
        \mk{nodes and arcs}.
  \begin{itemize}
    \item \textbf{Nodes} $=$ objects, concepts, situations or events.
    \item \textbf{Arcs} $=$ the relationship between two nodes, and the arc is
          \textbf{labelled with that relationship}.
  \end{itemize}
  \item A \textbf{declarative graphic representation}: it states what is true and supports
        automated reasoning over it, but contains no procedure.
  \item Proposed as \mk{an alternative to predicate logic}. Same facts,
        but all of one object's facts meet at one node instead of being scattered over
        clauses.
  \item \textbf{Oldest KR in the subject}: the Tree of Porphyry, 3rd century AD.
        Worth a line in an 8-mark answer.
\end{itemize}
\lead{The two links every paper wants named}
\begin{itemize}
  \item \code{isa} / \code{a-kind-of} $-$ \textbf{class to class}. \emph{A cat is a
        mammal}. Set inclusion: $Cats \subset Mammals$.
  \item \code{instance} / \code{member} $-$ \textbf{individual to class}. \emph{Tom is a
        cat}. Set membership: $Tom \in Cats$.
  \item \mk{Getting these two the wrong way round is the commonest lost
        mark} in the drawing questions. \emph{Sakti Gauchan} is an \code{instance} of
        Person; \emph{Person} is \code{isa} Mammal.
  \item Everything else is an ordinary relation arc: \code{has-part},
        \code{owns}, \code{likes}, \code{team}, \code{uniform-color}.
\end{itemize}
\figT{c5_isa_instance.png}
{\footnotesize\color{sub} \code{isa} runs class to class, \code{instance-of} runs
individual to class, and \code{owns} is an ordinary relation. Source: BA Sir CH-05.}}{%
\lead{The canonical net, and its FOPL}
\figT{c5_semnet_peewee.png}
{\footnotesize\color{sub} Rich and Knight's figure. \mk{Every
``convert these sentences'' question in the papers is this net with the names
changed}: a class chain at the top, an individual hung off it by
\code{instance}, and its own properties fanning out. Source: PS Sir ch5.}
\lead{Why a net can do things logic cannot}
\begin{itemize}
  \item \textbf{Inheritance is free.} A property put on \emph{Person} is automatically
        available to \emph{Pee-Wee-Reese} by following \code{instance} then \code{isa}
        upward. No inference rule is needed; the arcs \emph{are} the rule.
  \item \textbf{Intersection search.} Spread activation outward from two nodes and see
        where the two wavefronts meet. That is how a net answers
        \emph{``what connects the Brooklyn Dodgers and blue?''} $-$ a question logic can
        only answer by proving a conjecture you already had to guess.
  \item \mk{This is the advantage to quote}, because it exploits the
        entity-based organisation that slot-and-filler structures have and logic does not.
\end{itemize}}

\Q{Advantages and limitations of a semantic net; how do you represent a three-place relation?}{\m{2}\m{4}}
{\color{sub}\footnotesize part of the \emph{advantages and limitations} question in
\yr{73 Ma}, \yr{76 Ba}, \yr{\textbf{78 Ch}}, \yr{78 Ba}, \yr{\textbf{\texttt{81 Bh}}}
\gap$\cdot$\gap chip sits on \S5.6, where the comparison is}

\sbsr{0.50}{%
\lead{Advantages}
\begin{itemize}
  \item \textbf{Natural and transparent}: close to how people link ideas, so it is
        readable by a non-programmer.
  \item \textbf{Easy to visualise}, and easy to check by eye.
  \item \textbf{Inheritance} of properties comes for free down \code{isa} chains.
  \item \textbf{Space-efficient}: each object appears \emph{once}, however many facts
        mention it.
  \item Related facts sit next to each other, so \textbf{relevant knowledge is found by
        following links} instead of by searching the whole store.
\end{itemize}
\lead{Limitations}
\begin{itemize}
  \item \mk{Cannot express quantification, negation or disjunction.}
        There is no way to draw \emph{``every mammal has \emph{some} parent''},
        \emph{``no bird is a mammal''} or \emph{``Tom is black or white''}. This is the
        limitation to write first.
  \item \textbf{No standard meaning for link names.} \code{has-part} means whatever the
        person who drew it meant, so two nets are not comparable.
  \item \textbf{No inference mechanism of its own}: it stores, it does not deduce.
  \item \textbf{Slow at query time} $-$ answering may need a traversal of the whole graph.
  \item \textbf{Does not scale}: a net that really modelled human memory cannot be built
        by hand.
  \item Interpretation is \textbf{informal}, so the net depends entirely on its creator.
\end{itemize}}{%
\lead{Three-place relations: use an event node}
\begin{itemize}
  \item Arcs are \textbf{binary}, so \emph{``John gave the book to Mary''} cannot be one
        arc: it has three arguments.
  \item \mk{The fix}: create a \textbf{new node for the event itself},
        then hang one binary arc on it per argument.
\end{itemize}
\figT{c5_semnet_event.png}
{\footnotesize\color{sub} \code{EV7} is the giving-event node. \code{agent} $\rightarrow$
John, \code{object} $\rightarrow$ the book, \code{beneficiary} $\rightarrow$ Mary, and
\code{instance} says \code{EV7} is a \emph{Give}. Source: PS Sir ch5.}
\begin{itemize}
  \item Same trick for a baseball game with two teams and a score:
        \code{isa(G23, Game)}, \code{home\_team(G23, Dodgers)},
        \code{visiting\_team(G23, Cubs)}, \code{score(G23, 5-3)}.
  \item \mk{Say why it works}: it converts an $n$-place predicate into
        $n$ binary ones by \emph{reifying} the relation as an object.
\end{itemize}}

\Q{What do the arcs of a semantic net actually mean?}{\m{2}}
{\color{sub}\footnotesize not asked on its own \gap$\cdot$\gap included because it is the
\mk{one line that separates a 5-mark answer from a 7-mark one} on
every ``explain semantic nets'' question}

\sbsr{0.53}{%
\figT{c5_semnet_links.png}}{%
\begin{itemize}
  \item \code{Subset} arc $=$ \textbf{set inclusion}, $A \subset B$. Eg
        $Cats \subset Mammals$.
  \item \code{Member} arc $=$ \textbf{set membership}, $A \in B$. Eg $Felix \in Cats$.
  \item A named arc between two \emph{individuals} $=$ the ground fact $R(A,B)$. Eg
        \code{Friend(Felix, Opus)}.
  \item A named arc drawn out of a \emph{class} box means
        $\forall x\, (x \in A \Rightarrow R(x,B))$: \textbf{every} member has that value.
        Eg \code{Legs} out of Birds $\rightarrow 2$.
  \item Same arc with an existential reading:
        $\forall x\, (x \in A \Rightarrow \exists y\, (y \in B \wedge R(x,y)))$. Eg
        \code{Parent} out of Bird $\rightarrow$ Birds: every bird has \emph{some} bird
        parent.
  \item \mk{Why it matters}: the last two are the only quantification a
        net can express, and only because the reader agrees to read the picture that way.
        Nothing in the drawing says which of the two is meant. That ambiguity \emph{is}
        the ``lack of formality'' limitation.
  \item {\footnotesize\color{sub} Source: BJ Sir \code{AI\_Chapter5\_new}.}
\end{itemize}}

%-----------------------------------------------------------------------
\T{5.5 Frames}

\Q{What is a frame? Explain with a suitable example. What is its significance in AI?}{\m{1}\m{5}\m{7}\m{8}}
{\color{sub}\footnotesize \tF{5} \yr{69 Po \m{7}, \textbf{74 Bh} \m{7},
\textbf{\textit{71 Bh}} \m{8}, 82 Ka \m{1}, 80 Ash \m{5}} \gap$\cdot$\gap
\yr{82 Ka} wants \mk{one sentence} for 1 mark \gap$\cdot$\gap
\yr{80 Ash} asks for the \emph{importance} of frames, so end on the advantages}

\sbsr{0.50}{%
\lead{Definition}
\begin{itemize}
  \item A \textbf{frame} is a \mk{collection of attributes and their
        values that describes one entity}, together with that entity's relationship to
        other frames.
  \begin{itemize}
    \item \textbf{Slot} $=$ the attribute name. \textbf{Filler} $=$ its value.
    \item So a frame is \textbf{a record structure with meaning attached}.
  \end{itemize}
  \item \mk{A single frame is almost useless.} The power is in a
        \textbf{frame system}: many frames, connected because \textbf{the value of one
        frame's slot is another frame}.
  \item \mk{The line that earns the mark}: a semantic net is a
        \textbf{two-dimensional} representation, nodes and arcs;
        \textbf{frames add a third dimension by letting a node have internal structure}.
  \item A frame is designed to represent a \textbf{stereotypical} object or situation:
        what is normally true, with defaults for what has not been observed.
  \item Frames are \textbf{static}. A script (\S5.7) is the same idea made
        \emph{sequential}.
\end{itemize}
\lead{The four parts of a frame}
\begin{itemize}
  \item \textbf{1. Frame identification name} $-$ what entity this frame is about.
  \item \textbf{2. Relationship to other frames} $-$ \code{isa}, \code{instance},
        \code{part-of}, \code{superclass}.
  \item \textbf{3. Slot and filler} $-$ the attribute--value pairs themselves.
  \item \textbf{4. Default information} $-$ values taken as true \emph{until contradicted}.
        \mk{Logic has no clean way to do this}; frames do, and it is
        the strongest single argument for them.
\end{itemize}
\lead{Types of frame}
\begin{itemize}
  \item \textbf{Class frame} $-$ the generic frame others inherit from. Eg \emph{Vehicle}.
  \item \textbf{Subclass frame} $-$ inherits a class frame and adds slots. Eg \emph{Car}.
  \item \textbf{Instance frame} $-$ one individual. Eg \emph{Ram's Car}.
  \item An instance or subclass frame \textbf{inherits every slot above it} and may add
        its own. Generic frames sit at the top of the hierarchy.
\end{itemize}}{%
\lead{Class, subclass and instance in one picture}
\figT{c5_frame_classes.png}
{\footnotesize\color{sub} \emph{Vehicle} (class) $\rightarrow$ \emph{Truck} and
\emph{Car} (subclass, \code{is-a}) $\rightarrow$ \emph{Ram's Car} (instance), with
\code{part\_of} links out to \emph{Basket} and \emph{Ram}. Note the instance frame
\textbf{fills in} the slots it inherited empty. Source: \emph{Insights} book p136.}}

\penalty0\vspace{2pt}

\sbsr{0.41}{%
\figT{c5_frame_hotel.png}
{\footnotesize\color{sub} Part of a frame description of a hotel room. Source: BA Sir
CH-05.}}{%
\lead{A frame system: four frames pointing at each other}
\begin{itemize}
  \item \emph{hotel room} carries three ordinary slots (\emph{specialization of},
        \emph{location}, \emph{contains}), and the third one's filler is
        \mk{three pointers, not three strings}. That is what makes this a
        frame \emph{system}.
  \item Each frame it points at repeats the pattern. \emph{hotel chair} has
        \emph{specialization of: chair}, so it inherits everything a chair has and
        \textbf{overrides} only \emph{height}, \emph{legs} and \emph{use}.
  \item \emph{hotel bed} adds a \code{part} slot filled by the \emph{mattress} frame,
        which is \textbf{part-of}, not \textbf{is-a}. \mk{Confusing those
        two is the commonest error in a frame answer}: \emph{is-a} inherits,
        \emph{part-of} does not.
  \item \mk{The line to write under your own drawing}: the diagram is a
        graph whose nodes happen to be records. That is structured knowledge
        representation in one picture, and it is why \yr{\textbf{\texttt{80 Bh}}} can ask
        for a net and a frame in the same breath.
\end{itemize}}

\Q{What can a slot hold, and what are demons?}{\m{2}\m{4}}
{\color{sub}\footnotesize inside the frame question in \yr{\textbf{\textit{71 Bh}}}
\m{8} and \yr{69 Po} \m{7} \gap$\cdot$\gap
\mk{a slot holds more than a value}, and saying so is the difference
between a 5-mark frame answer and an 8-mark one}

\sbsr{0.42}{%
\figT{c5_frame_facets.png}
{\footnotesize\color{sub} One STUDENT frame with every facet filled in. Source: PS Sir
ch5.}
\figT{c5_frame_demons.png}
{\footnotesize\color{sub} A generic property frame carrying two demons:
\code{if-added: ADD\_PROPERTY} on \emph{types}, \code{if-needed: FIND\_OWNER} on
\emph{owner}, which also shows a \textbf{default} of \emph{government}. Source: PS Sir
ch5.}}{%
\lead{Facets: what a slot carries besides its value}
\begin{itemize}
  \item \textbf{Value} $-$ the current filler.
  \item \textbf{Default} $-$ the filler used when nothing has been observed.
  \item \textbf{Cardinality} $-$ minimum and maximum number of fillers allowed.
  \item \textbf{Type} $-$ a restriction on what a filler may be, usually another frame.
  \item \textbf{Procedures} $-$ attached code, see demons below.
  \item \textbf{Salience} $-$ how important this slot is to the frame.
  \item \textbf{Constraints} $-$ axioms the value must satisfy, eg $0 \le CGPA \le 4.0$.
\end{itemize}
\lead{\mk{Demons} (attached procedures, Insights \S5.4.3)}
\begin{itemize}
  \item A demon is a \textbf{procedure stored in a slot and fired automatically when the
        slot is touched}. The frame does not just hold data, it reacts.
  \item \textbf{IF-NEEDED} $-$ fires when a value is required and none is stored.
        \emph{Compute it now.}
  \item \textbf{IF-ADDED} $-$ fires when a value is put into the slot.
        \emph{Propagate the consequences.}
  \item \textbf{IF-REMOVED} $-$ fires when a value is deleted.
  \item \textbf{IF-CHANGED} $-$ fires when a stored value is replaced.
  \item \mk{Why this matters for the exam}: demons are what make a frame
        system \emph{active}, and they are the honest answer to
        \emph{``frames have no inference mechanism''} $-$ they have no \emph{general} one,
        but each slot can carry its own.
  \item \textbf{Frames Vs objects}: slots are data members, demons are methods, \code{isa}
        is inheritance and \code{part-of} is containment. \mk{A frame
        system is object-oriented programming, invented for AI a decade earlier.} One
        sentence, and it is often worth a mark.
\end{itemize}}

\Q{Advantages and limitations of frames}{\m{2}\m{4}}
{\color{sub}\footnotesize part of the \emph{advantages and limitations} question,
chip on \S5.6 \gap$\cdot$\gap \mk{write these as a pair with the semantic
net list}, because every paper that asks for one asks for the other}

\sbs{%
\lead{Advantages}
\begin{itemize}
  \item \textbf{Groups related knowledge together}, so programming and retrieval are
        easier: everything about an object is in one record.
  \item \textbf{Defaults}, and therefore the ability to reason with incomplete
        information and to detect missing values.
  \item \textbf{Easy to extend}: adding a slot, a property or a relation costs one line.
  \item \textbf{Inheritance} down the class hierarchy, exactly as in a semantic net.
  \item \textbf{Procedural attachment}: a slot can compute or validate itself.
  \item \textbf{Flexible and widely used} $-$ vision and natural language processing both
        run on frames.
  \item Good at \textbf{causal knowledge}, because the information is organised by cause
        and effect.
\end{itemize}}{%
\lead{Limitations}
\begin{itemize}
  \item \mk{No inference mechanism of its own.} Frames store and
        inherit; they do not deduce. Any reasoning has to be bolted on.
  \item \textbf{Modification is dangerous}: change one frame and every frame that inherits
        from it is affected.
  \item \textbf{Redundancy} $-$ the same fact easily ends up in several frames.
  \item \textbf{No standard structure}: the frame for a room and the frame for a person
        look nothing alike, so nothing generic can be written over them.
  \item \textbf{Defaults can be wrong}, and nothing in the formalism flags it when they
        are.
  \item \mk{The examiner's expected trade-off line}: frames win on
        organisation and defaults, logic wins on proof, nets win on visibility.
\end{itemize}}

%-----------------------------------------------------------------------
\T{5.6 Semantic nets Vs frames, and converting one into the other}

\Q{Explain semantic nets and frames with examples. Compare them, and list their advantages and limitations}{\m{2}\m{4+4}\m{6+2}\m{7}}
{\color{sub}\footnotesize \tS{13} \yr{\textbf{72 Ash} \m{7}, 72 Ma \m{7}, 73 Ma \m{4+4},
\textit{74 Ma} \m{2+5}, \textbf{74 Bh} \m{7}, \textbf{\textit{74 Bh}} \m{4+3},
75 Ba \m{2+6}, 76 Ba \m{6+2}, 78 Ba \m{6+2}, \textbf{78 Ch} \m{4+4}, 79 Jth \m{5},
\texttt{81 Ba} \m{3+5}, \textbf{\texttt{81 Bh}} \m{2}}
\gap$\cdot$\gap \mk{the second-most-asked question in the chapter}
\gap$\cdot$\gap \yr{\textbf{72 Ash}} and \yr{78 Ba} ask the
\mk{``why are they important in AI''} follow-up, so keep two lines for it}

\begin{tabularx}{\linewidth}{@{}L{2.5cm}L{5.0cm}X@{}}
\toprule
\hrow & \thead{Semantic net} & \thead{Frame} \\
\midrule
Shape & labelled directed \textbf{graph} & collection of \textbf{records}, slots and
fillers \\
Dimension & \textbf{two}: node and arc & \textbf{three}: node, and structure
\emph{inside} the node \\
Unit of knowledge & one \textbf{arc} $=$ one fact & one \textbf{frame} $=$ one whole
object \\
Node holds & a name only & slots, values, defaults, types, constraints, procedures \\
Relations & drawn as \textbf{arcs} & written as \textbf{slot values that are other
frames} \\
Inheritance & \checkmark down \code{isa} / \code{instance} arcs &
\checkmark down class $\rightarrow$ subclass $\rightarrow$ instance \\
Defaults & \xmark cannot express & \checkmark first-class, with
\code{if-needed} to compute one \\
Procedures & \xmark none & \checkmark demons on any slot \\
Best at & \textbf{seeing} the relationships at a glance & \textbf{holding} everything
about one object in one place \\
Weakest at & complex or structured facts about a single node & showing the shape of the
knowledge as a whole \\
Reasoning & intersection search, link following & inheritance and slot lookup only \\
Both fail at & \multicolumn{2}{@{}X@{}}{quantification, negation, disjunction, and having
no inference engine of their own} \\
\bottomrule
\end{tabularx}
\vspace{2pt}

\sbs{%
\lead{\mk{How frames are useful in semantic nets}
\yr{\textbf{72 Ash}} \m{7}}
\begin{itemize}
  \item The question sounds like two topics. It is one: \textbf{a frame is what you get
        when you let a semantic-net node have contents}.
  \item A net node called \emph{Tom} can only be named. Make it a frame and it can carry
        \emph{colour: ginger}, \emph{age: default 2}, \emph{owner: John}, plus a demon
        that computes age from a birth date.
  \item So frames \textbf{fix the net's weakest point}: they let one node hold structured,
        defaulted, constrained knowledge instead of forcing every detail out onto its own
        arc.
  \item And the net \textbf{fixes the frame system's weakest point}: it shows the shape of
        the whole knowledge base, which a pile of records does not.
  \item \mk{Answer in one sentence}: frames give a semantic net depth at
        each node, and the net gives a frame system its map.
\end{itemize}
\lead{Why both matter in AI \yr{78 Ba} \m{2}}
\begin{itemize}
  \item They give \textbf{inheritance and defaults}, which is how an agent reasons with
        incomplete knowledge.
  \item They are \textbf{human-readable}, so a domain expert can check the knowledge base
        without knowing logic $-$ the bottleneck in every expert system (Ch\,7).
  \item They organise knowledge by \textbf{object} rather than by \emph{sentence}, so
        retrieval is local instead of a search of the whole store.
\end{itemize}}{%
\lead{\mk{Converting a semantic net into a frame}
\yr{\texttt{81 Ba}} \m{5}}
\begin{itemize}
  \item \textbf{Step 1.} Every \textbf{node} becomes a \textbf{frame}, named after the
        node.
  \item \textbf{Step 2.} Every \textbf{outgoing arc} of that node becomes a \textbf{slot}
        in its frame. \mk{The arc label is the slot name, the arc's
        target is the filler.}
  \item \textbf{Step 3.} \code{isa} and \code{instance} arcs become the frame's
        \textbf{class link}, not ordinary slots: \code{isa} $\rightarrow$ a subclass
        frame, \code{instance} $\rightarrow$ an instance frame.
  \item \textbf{Step 4.} Any slot whose filler is a node that has its own frame is
        \textbf{left as a pointer to that frame}, which is what makes the result a frame
        \emph{system}.
  \item \textbf{Step 5.} Push each property to the \textbf{highest} class that has it, and
        let everything below inherit. A property drawn on three individuals in the net
        belongs on their class in the frame system.
  \item \mk{Check}: no information may be lost. Count the arcs in the
        net; count the filled slots in the frames. They must match.
  \item \yr{\texttt{81 Ba}} is worked end to end in the companion, on this net.
\end{itemize}}

\sbsr{0.42}{%
\lead{The same knowledge three ways}
\begin{itemize}
  \item \mk{Frame, net and FOPL are interchangeable} for taxonomic
        knowledge, and \yr{75 Ba} \m{6} asks for exactly that: \emph{examples of both,
        with FOPL statements}.
  \item Read the figure left to right: a frame-based knowledge base, then the
        first-order sentences that say the same thing.
  \item \code{Subset(Birds, Animals)} is the \code{isa} arc.
        \code{Member(Opus, Penguins)} is the \code{instance} arc.
        \code{Rel(Legs, Birds, 2)} is a slot with a value.
  \item \mk{Write the FOPL line beside every arc you draw} in that
        question. It is 2 free marks and most answers skip it.
\end{itemize}}{%
\figT{c5_frame_fopl.png}
{\footnotesize\color{sub} (a) a frame-based knowledge base and (b) its translation into
first-order logic, the same nine facts twice. Source: PS Sir ch5.}}

%-----------------------------------------------------------------------
\T{5.7 Conceptual dependency and scripts}

\Q{What is conceptual dependency? Explain the common primitives used in it}{\m{2+5}\m{3}}
{\color{sub}\footnotesize \tP{2} \yr{\textbf{71 Bh} \m{2+5}, \textbf{77 Ch} \m{3+5}}
\gap$\cdot$\gap only 2 papers, but \mk{both times CD and scripts
together carry the whole 8-mark question} \gap$\cdot$\gap
\yr{\textbf{77 Ch}} pairs it with \emph{``how is knowledge represented using scripts''}}

\sbsr{0.46}{%
\lead{What it is}
\begin{itemize}
  \item \textbf{Schank, 1973--75}, built for \textbf{natural language understanding}.
  \item A representation of the \mk{meaning} of a sentence, built from
        a \textbf{small fixed set of primitive acts} instead of from the sentence's own
        words.
  \item \textbf{Three goals}, and the papers want all three:
  \begin{itemize}
    \item to \textbf{help draw inferences} from a sentence;
    \item to be \textbf{independent of the words} used in the input;
    \item so that \mk{any two sentences identical in meaning have
          exactly one representation}.
  \end{itemize}
  \item Hence \emph{I gave the man a book}, \emph{the man got a book from me} and
        \emph{the book was given to the man by me} all draw the \textbf{same} diagram.
        That is the point of the whole theory.
  \item A \textbf{strong} slot-and-filler structure: unlike a frame, the slots have
        \textbf{fixed, agreed meanings}.
\end{itemize}
\lead{The four conceptual categories}
\begin{itemize}
  \item \textbf{ACT} $-$ an action, always one of the primitives below.
  \item \textbf{PP} $-$ \emph{picture producer}, an object or actor.
  \item \textbf{AA} $-$ \emph{action aider}, a modifier of an ACT (how, when, where).
  \item \textbf{PA} $-$ \emph{picture aider}, a modifier of a PP (colour, size, quality).
  \item \textbf{Dependency rule}: PP and ACT are \textbf{governors}, PA and AA are
        \textbf{dependents}. A dependent cannot stand alone; a conceptualization needs at
        least two governors.
  \item Eg \emph{Sally stroked her fat cat}: PP $=$ Sally, cat; ACT $=$ stroke;
        PA $=$ fat.
\end{itemize}}{%
\lead{The eleven primitive acts}
\begin{tabularx}{\linewidth}{@{}L{1.85cm}X@{}}
\toprule
\hrow \thead{Primitive} & \thead{Transfer or action, and a verb that maps to it} \\
\midrule
\textbf{ATRANS} & transfer of an \textbf{abstract relationship}, ie possession. \emph{give,
lend, borrow} \\
\textbf{PTRANS} & transfer of the \textbf{physical location} of an object. \emph{go, come,
went} \\
\textbf{PROPEL} & application of a \textbf{physical force} to an object. \emph{push, pull}
\\
\textbf{MOVE} & movement of a \textbf{body part} by its owner. \emph{kick, punch} \\
\textbf{GRASP} & an actor \textbf{grasping} an object. \emph{clutch, catch} \\
\textbf{INGEST} & an animal \textbf{taking something in}. \emph{eat, drink} \\
\textbf{EXPEL} & an animal \textbf{forcing something out}. \emph{cry, sweat, throw} \\
\textbf{MTRANS} & transfer of \textbf{mental information}. \emph{tell, ask, read} \\
\textbf{MBUILD} & \textbf{building new information} out of old. \emph{decide, infer,
conclude} \\
\textbf{SPEAK} & \textbf{producing a sound}. \emph{say, talk, play music} \\
\textbf{ATTEND} & \textbf{focusing a sense organ} on a stimulus. \emph{listen, watch} \\
\bottomrule
\end{tabularx}
\vspace{2pt}
\begin{itemize}
  \item \mk{Learn the first three and MTRANS.} Every example in every
        paper and every scene of the restaurant script is built from those four.
\end{itemize}}

\Q{Draw the conceptual dependency of a sentence}{\m{2}\m{5}}
{\color{sub}\footnotesize inside \yr{\textbf{71 Bh}} \m{5} \gap$\cdot$\gap
the notation cannot be typed, so \mk{learn it as a picture}}

\sbsr{0.58}{%
\figT{c5_cd_examples.png}
{\footnotesize\color{sub} The seven dependency forms, then each one used on a sentence.
Row 6 is \emph{John pushed the cart} ($ACT \xleftarrow{o} PP$, PROPEL); row 7 is
\emph{John took the book from Mary} (ATRANS with a recipient pair). Source: BA Sir
CH-05.}}{%
\lead{The notation}
\begin{itemize}
  \item $\Leftrightarrow$ \textbf{two-way link} between an actor and the action. Every
        conceptualization has exactly one.
  \item $\rightarrow$ \textbf{direction of dependency}: the dependent points at its
        governor.
  \item \textbf{o} $=$ objective case, what the act is done \emph{to}.
  \item \textbf{R} $=$ recipient case, drawn as a \textbf{fork}: the arrow up is the
        receiver, the arrow down is the giver.
  \item \textbf{I} $=$ instrumental case, what it was done \emph{with}.
  \item \textbf{D} $=$ directive case, from where to where.
  \item \textbf{p} $=$ past tense, \textbf{f} future, \textbf{t} transition,
        \textbf{?} interrogative, \textbf{/} negative. Tense sits on the
        $\Leftrightarrow$.
  \item \mk{Worked example} $-$ \emph{I gave a book to Ram}:
  \begin{itemize}
    \item $I \Leftrightarrow ATRANS$, with \textbf{p} on the link;
    \item $\xleftarrow{o}$ book;
    \item $\xleftarrow{R}$ forking up to \emph{Ram} and down to \emph{I}.
  \end{itemize}
\end{itemize}
\lead{Problems with CD}
\begin{itemize}
  \item \textbf{The original sentence cannot be reconstructed} from the diagram; meaning
        survives, wording does not.
  \item \textbf{Choosing the primitive is hard}: one verb may map to several ACTs
        depending on context.
  \item \textbf{The structures get huge.} \emph{``John bet Mike that Nepal will win the
        World Cup''} needs an enormous diagram for one sentence.
  \item Rules must be \textbf{hand-written per primitive}, so the theory does not scale to
        a general-purpose KR.
\end{itemize}}

\Q{How is knowledge represented using scripts? Explain with an example}{\m{5}}
{\color{sub}\footnotesize \yr{\textbf{77 Ch} \m{5}} \gap$\cdot$\gap 1 paper at 5 marks,
so no chip of its own \gap$\cdot$\gap \mk{always asked as the second half of
the CD question}, never alone \gap$\cdot$\gap the restaurant script is the only example
any deck or the textbook uses, so learn that one}

\sbsr{0.47}{%
\lead{What a script is}
\begin{itemize}
  \item \textbf{Schank and Abelson, 1977}, written in the CD notation.
  \item A \mk{structured description of a stereotyped sequence of
        events in a known context}: going to a restaurant, going to the theatre, visiting
        a doctor.
  \item It \textbf{extends a frame} by making the slot values \textbf{ordered in time} and
        by representing \emph{expectations} about what happens next.
  \item Justified by three facts about the world: events occur in known patterns, events
        cause one another, and entry conditions decide whether a pattern applies at all.
\end{itemize}
\lead{\mk{The six components} (write all six)}
\begin{itemize}
  \item \textbf{Entry conditions} $-$ must be true before the script can be invoked.
        Eg \emph{S is hungry}, \emph{S has money}.
  \item \textbf{Results} $-$ true once it has run. Eg \emph{S is not hungry},
        \emph{O has more money}.
  \item \textbf{Props} $-$ the objects involved. Eg tables, menu, food, check, money.
  \item \textbf{Roles} $-$ the participants and what each does. Eg S $=$ customer,
        W $=$ waiter, C $=$ cook, M $=$ cashier, O $=$ owner.
  \item \textbf{Track} $-$ a variation on the general pattern. Eg \emph{coffee shop} Vs
        \emph{fancy restaurant}: same script, different track.
  \item \textbf{Scenes} $-$ the ordered sequence of events, each written in CD primitives.
\end{itemize}
\lead{The RESTAURANT script, header and scene 1}
\begin{itemize}
  \item \textbf{Script}: RESTAURANT. \textbf{Track}: coffee shop.
  \item \textbf{Props}: tables, menu, F $=$ food, check, money.
  \item \textbf{Roles}: S $=$ customer, W $=$ waiter, C $=$ cook, M $=$ cashier,
        O $=$ owner.
  \item \textbf{Entry conditions}: S is hungry, S has money.
  \item \textbf{Results}: S has less money, O has more money, S is not hungry,
        S is pleased (optional).
  \item \textbf{Scene 1, entering}: S PTRANS S into restaurant $\cdot$ S ATTEND eyes to
        tables $\cdot$ S MBUILD where to sit $\cdot$ S PTRANS S to table $\cdot$ S MOVE S
        to sitting position.
\end{itemize}}{%
\figT{c5_script_rest1.png}
{\footnotesize\color{sub} Scene 2, ordering, with both tracks: menu already on the table,
or the waiter brings it. Source: BA Sir CH-05.}}

\penalty0\vspace{2pt}

\sbsr{0.47}{%
\lead{Invocation, and why it is the hard part}
\begin{itemize}
  \item A script is opened when its topic is \textbf{significant} in the input; if merely
        mentioned, only a pointer is kept.
  \item \emph{``John enjoyed the play in the theatre''} $\rightarrow$ invoke, and every
        implicit question is answered: did he go, did he buy a ticket, did he have money.
  \item \mk{The failure case to quote}: \emph{``John went to the theatre
        to pick up his daughter''}. Same words, low significance, and invoking the script
        produces confident wrong answers.
  \item There is no reliable rule for judging significance, only heuristics. That is the
        honest limitation to close on.
\end{itemize}
\lead{Advantages and disadvantages}
\begin{itemize}
  \item \checkmark \textbf{Predicts implicit events}: it fills in what
        the text never said, which is exactly what language understanding needs.
  \item \checkmark Builds \textbf{one coherent interpretation} out of
        scattered observations.
  \item \xmark \mk{Far more specific and less
        general than a frame.} One script covers one situation.
  \item \xmark Not suitable for knowledge that is not
        event-sequenced.
  \item \xmark Wrong invocation gives confidently wrong
        answers.
  \item The cure for the inflexibility: \textbf{MOPs}, memory organization packets $-$
        smaller reusable modules recombined to fit the situation.
\end{itemize}}{%
\figT{c5_script_rest2.png}
{\footnotesize\color{sub} Scene 3 eating and scene 4 exiting, including the
\emph{no pay} path. Source: BA Sir CH-05.}}

\hr

{\color{sub}\footnotesize \textbf{Source notes.}
\emph{Insights} \S5.4.2 (book p135) converts the Tom-the-cat semantic net into frames, but
its figure prints \code{like John} and \code{type Ginger} where the net says
\emph{is-owned-by} and \emph{is-coloured}, and gives the \emph{John} frame an
\code{Age 22 / Color white} pair carried over from the car example on the facing page.
The net-to-frame recipe in \S5.6 above follows the book's \emph{method} and Rich and
Knight's \emph{figures}, and the companion works the conversion on a net that is drawn
correctly.
\gap$\cdot$\gap
The syllabus files this chapter as five sub-topics but the papers only ever set four:
representations and mappings is never asked on its own, always as the two-mark opener to
\emph{what is knowledge representation}.}
